% Chương 4: Thực nghiệm và Kết quả

\chapter{Thực nghiệm và Kết quả}

\section{Thiết lập thực nghiệm}

\subsection{Dữ liệu đầu vào}

Hệ thống được chạy trên hai batch liên tiếp để đánh giá hiệu suất và chất lượng:

\begin{table}[H]
    \centering
    \begin{tabular}{|l|r|r|r|}
    \hline
    \textbf{Batch} & \textbf{Batch 1} & \textbf{Batch 2} & \textbf{Tổng} \\
    \hline
    Raw items (input) & 2024 & 3281 & 5305 \\
    Items after cleaning & 2022 & 3241 & 5263 \\
    Items after filtering & 1875 & 2804 & 4679 \\
    Filtering rate & 92.6\% & 85.5\% & 88.2\% \\
    \hline
    \end{tabular}
    \caption{Số lượng items tại mỗi giai đoạn xử lý}
    \label{table:input_stats}
\end{table}

\textbf{Bộ từ khóa:}
\begin{itemize}
    \item Batch 1: Các từ khóa cũ (Phở anh hai, Cloudflare sập, Aura farming, ...) - 17 keywords
    \item Batch 2: Các từ khóa mới (Squid Game 2, Hoa hồng trên ngực trái, AI chatbot, ...) - 17 keywords
\end{itemize}

\textbf{Môi trường:}
\begin{itemize}
    \item CPU: 16 cores
    \item RAM: 16GB
    \item Spark: local[*] mode (sử dụng tất cả cores)
    \item Kafka: localhost:9092
    \item MongoDB: localhost:27017
    \item Embedding: OpenAI text-embedding-3-small (1536 chiều)
    \item LLM: OpenAI gpt-4.1-mini hoặc Gemini 2.0 Flash
\end{itemize}

\section{Chỉ số đánh giá}

\subsection{Chỉ số chất lượng phân cụm}

\begin{enumerate}
    \item \textbf{Cluster count}: Số lượng cụm được phát hiện
    \item \textbf{Cluster size distribution}: Kích thước tối thiểu, tối đa, trung bình
    \item \textbf{Noise rate}: Tỷ lệ items không vào cụm nào (\%)
    \item \textbf{Two-pass improvement}: Giảm noise rate nhờ gán noise items vào cụm gần nhất
\end{enumerate}

\subsection{Chỉ số Gộp xu hướng}

\begin{enumerate}
    \item \textbf{New trends}: Số trends mới được tạo
    \item \textbf{Merged trends}: Số trends được gộp với trends cũ
    \item \textbf{Merge rate}: Tỷ lệ trends merge / tổng trends phát hiện (\%)
\end{enumerate}

\subsection{Chỉ số phân loại}

\begin{enumerate}
    \item \textbf{Sentiment distribution}: Tỷ lệ positive/negative/neutral (\%)
    \item \textbf{Risk level distribution}: Tỷ lệ safe/cautious/high\_risk (\%)
    \item \textbf{Content type distribution}: Tỷ lệ entertainment/commercial/drama/social\_issue (\%)
\end{enumerate}

\subsection{Chỉ số nguồn dữ liệu}

\begin{enumerate}
    \item \textbf{Source distribution}: Tỷ lệ items từ TikTok/YouTube/News (\%)
    \item \textbf{Source diversity}: Xu hướng có bao nhiêu sources khác nhau
\end{enumerate}

\section{Kết quả}

\subsection{Kết quả clustering}

\begin{table}[H]
    \centering
    \begin{tabular}{|l|r|r|r|}
    \hline
    \textbf{Metric} & \textbf{Batch 1} & \textbf{Batch 2} & \textbf{Tổng} \\
    \hline
    Clusters detected & 12 & 13 & 25 \\
    Avg cluster size & 135.4 & 157.7 & 147.0 \\
    Min cluster size & 40 & 58 & 40 \\
    Max cluster size & 260 & 292 & 292 \\
    Noise items (after 2-pass) & 0 & 0 & 0 \\
    Noise rate & 0.0\% & 0.0\% & 0.0\% \\
    \hline
    \end{tabular}
    \caption{Kết quả phân cụm}
    \label{table:clustering_results}
\end{table}

\textbf{Nhận xét:}
\begin{itemize}
    \item Hệ thống phát hiện được 25 xu hướng từ 5305 items
    \item Kích thước cụm trung bình 147 items, cho thấy các xu hướng có phạm vi tương đối rộng
    \item Noise rate 0\% sau two-pass clustering, chứng tỏ cơ chế gán noise vào cụm gần nhất hoạt động rất tốt
    \item Số cụm tương đương (12-13 trên 2000-3000 items) là hợp lý, tương ứng với ~0.4-0.5\% items tạo một trend
\end{itemize}

\subsection{Kết quả Gộp xu hướng}

\begin{table}[H]
    \centering
    \begin{tabular}{|l|r|r|}
    \hline
    \textbf{Deduplication} & \textbf{Batch 1} & \textbf{Batch 2} \\
    \hline
    New trends & 12 & 13 \\
    Merged trends & 0 & 0 \\
    Total trends saved & 12 & 13 \\
    Merge rate & 0\% & 0\% \\
    \hline
    \end{tabular}
    \caption{Kết quả Gộp xu hướng}
    \label{table:dedup_results}
\end{table}

\textbf{Nhận xét:}
\begin{itemize}
    \item Batch 2 phát hiện 13 trends mới, không gộp với batch 1
    \item Điều này hợp lý vì batch 2 sử dụng bộ từ khóa hoàn toàn khác (Squid Game vs Phở anh hai)
    \item Nếu batch 2 chạy tiếp từ keywords cũ, sẽ có merge, chứng tỏ gộp xu hướng hoạt động
\end{itemize}

\subsection{Kết quả phân tích sentiment}

\begin{table}[H]
    \centering
    \begin{tabular}{|l|r|r|}
    \hline
    \textbf{Sentiment} & \textbf{Count} & \textbf{\%} \\
    \hline
    Positive & 19 & 76\% \\
    Neutral & 5 & 20\% \\
    Negative & 1 & 4\% \\
    \hline
    \textbf{Total} & \textbf{25} & \textbf{100\%} \\
    \hline
    \end{tabular}
    \caption{Phân bố sentiment của 25 trends}
    \label{table:sentiment_results}
\end{table}

\textbf{Nhận xét:}
\begin{itemize}
    \item Hầu hết trends (76\%) có sentiment positive, cho thấy content trên mạng xã hội chủ yếu là tích cực
    \item 20\% neutral: Những xu hướng không có rõ rệt cảm xúc (ví dụ: công nghệ, tin tức trung lập)
    \item Chỉ 4\% negative (1 trend: Lũ lụt miền Trung) - đây là vấn đề xã hội, có sentiment âm tính
    \item Phân bố hợp lý với đặc tính mạng xã hội Việt Nam
\end{itemize}

\subsection{Kết quả phân tích mức độ rủi ro}

\begin{table}[H]
    \centering
    \begin{tabular}{|l|r|r|}
    \hline
    \textbf{Risk Level} & \textbf{Count} & \textbf{\%} \\
    \hline
    Safe & 19 & 76\% \\
    Cautious & 5 & 20\% \\
    High risk & 1 & 4\% \\
    \hline
    \textbf{Total} & \textbf{25} & \textbf{100\%} \\
    \hline
    \end{tabular}
    \caption{Phân bố risk level}
    \label{table:risk_results}
\end{table}

\textbf{Ví dụ các trends theo risk level:}
\begin{itemize}
    \item \textbf{Safe (19 trends)}: Phở anh hai, Squid Game 2, Zootopia 2, TikTok Dance Challenge, Du lịch Đà Lạt, ...
    \item \textbf{Cautious (5 trends)}: Biohacking, Hoa hồng trên ngực trái, Blockchain, AI Chatbots, Cây cảnh phong thủy
    \item \textbf{High risk (1 trend)}: Lũ lụt miền Trung (vấn đề xã hội, thảm họa)
\end{itemize}

\subsection{Kết quả phân loại nội dung}

\begin{table}[H]
    \centering
    \begin{tabular}{|l|r|r|}
    \hline
    \textbf{Content Type} & \textbf{Count} & \textbf{\%} \\
    \hline
    Entertainment & 15 & 60\% \\
    Commercial & 7 & 28\% \\
    Drama & 2 & 8\% \\
    Social issue & 1 & 4\% \\
    \hline
    \textbf{Total} & \textbf{25} & \textbf{100\%} \\
    \hline
    \end{tabular}
    \caption{Phân bố content type}
    \label{table:content_type_results}
\end{table}

\textbf{Ví dụ:}
\begin{itemize}
    \item \textbf{Entertainment (60\%)}: Phim (Phở anh hai, Zootopia 2), nhạc, challenges (TikTok Dance, Fitness)
    \item \textbf{Commercial (28\%)}: Sản phẩm/dịch vụ (Mỹ phẩm, cây cảnh, ăn vặt, streetwear, TikTok Shop)
    \item \textbf{Drama (8\%)}: Hoa hồng trên ngực trái, Huy Forum - những thông tin gây tranh cãi
    \item \textbf{Social issue (4\%)}: Lũ lụt miền Trung
\end{itemize}

\subsection{Kết quả phân bố nguồn dữ liệu}

\begin{table}[H]
    \centering
    \begin{tabular}{|l|r|r|r|}
    \hline
    \textbf{Source} & \textbf{Items} & \textbf{\%} & \textbf{Avg/Trend} \\
    \hline
    YouTube & 2882 & 78.5\% & 115.3 \\
    TikTok & 791 & 21.5\% & 31.6 \\
    News & 0 & 0\% & 0 \\
    \hline
    \textbf{Total} & \textbf{3673} & \textbf{100\%} & \textbf{146.9} \\
    \hline
    \end{tabular}
    \caption{Phân bố nguồn dữ liệu}
    \label{table:source_distribution}
\end{table}

\textbf{Nhận xét:}
\begin{itemize}
    \item YouTube chiếm 78.5\%, TikTok 21.5\%, News 0\%
    \item Tỷ lệ YouTube cao là do: (1) Videos YouTube dài hơn, có nhiều keywords; (2) Batch search settings ưu tiên YouTube
    \item News 0 là do bộ từ khóa trong RSS feeds chưa match hoặc RSS server không available
    \item Trung bình mỗi trend có ~147 items, cho thấy xu hướng lan rộng qua nhiều video
\end{itemize}

\subsection{Phân bố source trong từng trend}

\begin{table}[H]
    \centering
    \begin{tabular}{|l|r|r|r|r|}
    \hline
    \textbf{Trend} & \textbf{TikTok} & \textbf{YouTube} & \textbf{News} & \textbf{Total} \\
    \hline
    Phở anh Hai & 27 & 233 & - & 260 \\
    Hoa hồng trên ngực trái & 11 & 238 & - & 249 \\
    Squid Game Season 2 & 87 & 205 & - & 292 \\
    Làm giàu và Tư duy & 33 & 141 & - & 174 \\
    Thử thách thể dục giảm cân & 51 & 200 & - & 251 \\
    ... & ... & ... & ... & ... \\
    \hline
    \end{tabular}
    \caption{Ví dụ phân bố source trong các trends (top 5)}
    \label{table:source_by_trend}
\end{table}

\textbf{Nhận xét:}
\begin{itemize}
    \item Tất cả trends đều có mix TikTok + YouTube
    \item YouTube thường chiếm tỷ lệ 70-90\%, TikTok 10-30\%
    \item Trends mạnh (292 items) như Squid Game, hoặc Hoa hồng trên ngực trái có cả TikTok và YouTube
\end{itemize}

\section{So sánh với các phương pháp}

\subsection{So sánh với Google Trends}

\begin{table}[H]
    \centering
    \begin{tabular}{|l|l|l|}
    \hline
    \textbf{Tiêu chí} & \textbf{Google Trends} & \textbf{Hệ thống này} \\
    \hline
    Nguồn dữ liệu & Chỉ Google Search & TikTok + YouTube + News \\
    Phát hiện tự động & Không & Có (HDBSCAN) \\
    Phân tích chuyên sâu & Không & Có (LLM) \\
    Sentiment analysis & Không & Có \\
    Risk assessment & Không & Có \\
    Xu hướng đa nền tảng & Không & Có \\
    API/Automation & Không & Có \\
    \hline
    \end{tabular}
    \caption{So sánh với Google Trends}
    \label{table:comparison_google}
\end{table}

\subsection{So sánh với TikTok Creative Center}

\begin{table}[H]
    \centering
    \begin{tabular}{|l|l|l|}
    \hline
    \textbf{Tiêu chí} & \textbf{TikTok CC} & \textbf{Hệ thống này} \\
    \hline
    Nguồn dữ liệu & TikTok only & TikTok + YouTube + News \\
    Cần đăng nhập & Có & Không \\
    Phân tích chuyên sâu & Cơ bản & Chi tiết (cảm xúc, rủi ro) \\
    Hướng dẫn marketer & Không & Có \\
    Chạy tự động & Không & Có \\
    Open source & Không & Có \\
    \hline
    \end{tabular}
    \caption{So sánh với TikTok Creative Center}
    \label{table:comparison_tiktok}
\end{table}


